\documentclass[12pt, a4paper]{article}

\usepackage[utf8]{inputenc}
\usepackage[T1]{fontenc}
\usepackage{geometry}
\usepackage{setspace}
\usepackage{amsmath}
\usepackage{xcolor}
\usepackage{newtxtext,newtxmath}
\usepackage{enumitem}
\usepackage{fancyhdr}
\usepackage{hyperref}
\usepackage{titlesec}
\usepackage{graphicx}
\usepackage{multicol}
\usepackage{booktabs}
\usepackage{array}
\usepackage{longtable}

\definecolor{NTRA_Blue}{RGB}{0,82,155}
\definecolor{correct}{RGB}{16,185,129}

\geometry{a4paper, margin=1in, headheight=14pt}
\setstretch{1.15}
\setlength{\parindent}{0pt}
\setlength{\parskip}{0.3em}

\hypersetup{colorlinks=true, linkcolor=NTRA_Blue, urlcolor=NTRA_Blue}

\titleformat{\section}{\Large\bfseries\color{NTRA_Blue}}{\thesection}{0.8em}{}[\titlerule]
\titleformat{\subsection}{\large\bfseries}{\thesubsection}{0.8em}{}

\pagestyle{fancy}
\fancyhf{}
\fancyhead[L]{LinearMind MCQ Questions}
\fancyhead[R]{\nouppercase{\leftmark}}
\fancyfoot[C]{\thepage}
\renewcommand{\headrulewidth}{0.4pt}

\newcommand{\question}[1]{\vspace{0.6em}\noindent\textbf{#1}}

\begin{document}

\begin{titlepage}
    \centering
    \thispagestyle{empty}

    \begin{minipage}[c]{0.2\textwidth}
        \begin{flushleft}
            \includegraphics[width=2.5cm]{Linear-Depression/linearmind/public/illustrations/Picture1.png}
        \end{flushleft}
    \end{minipage}%
    \hfill
    \begin{minipage}[c]{0.5\textwidth}
        \centering
        \setstretch{1.1}
        {\large \textbf{Cairo University}} \\
        {\large Faculty of Engineering} \\
        {\large Electronics \& Communications Dept.}
    \end{minipage}%
    \hfill
    \begin{minipage}[c]{0.2\textwidth}
        \begin{flushright}
            \includegraphics[width=2.5cm]{Linear-Depression/linearmind/public/illustrations/Picture2.jpg}
        \end{flushright}
    \end{minipage}

    \vspace{3cm}
    \rule{\linewidth}{0.8mm} \\[0.5cm]
    {\color{NTRA_Blue} \huge \bfseries MCQ Questions \& Answers} \\[0.3cm]
    {\color{NTRA_Blue} \LARGE LinearMind Platform} \\[0.4cm]
    {\color{NTRA_Blue} \large 82 Questions across 12 Modules} \\
    \rule{\linewidth}{0.8mm}

    \vfill
    \begin{center}
        \setstretch{1.5}
        \textbf{Prepared by:} \\[0.3cm]
        \begin{tabular}{l r}
            \textbf{Abdallah Essam Mohamed Mahmoud} & \hspace{1cm} ID: 9220479 \\
            \textbf{Zeyad Antar Othman Abu-Zaid} & \hspace{1cm} ID: 9220331 \\
            \textbf{Abdelrahman Hatem Nasr Youssef} & \hspace{1cm} ID: 9220461 \\
        \end{tabular}
    \end{center}
    \vfill
    \rule{0.3\linewidth}{0.2mm} \\
    {\large May 2026}
\end{titlepage}

\tableofcontents
\newpage

%% ============================================================
%%                      QUESTIONS
%% ============================================================

\section{Introduction (6 Questions)}

\question{Q1. What does Linear Regression try to find?}
\begin{enumerate}[label=\alph*)]
    \item The average of all data points
    \item The best straight line that fits the data
    \item The maximum value in the dataset
    \item Clusters in the data
\end{enumerate}

\question{Q2. Which type of problem is Linear Regression designed to solve?}
\begin{enumerate}[label=\alph*)]
    \item Classification: assigning categories
    \item Regression: predicting continuous values
    \item Clustering: grouping similar items
    \item Ranking: ordering items
\end{enumerate}

\question{Q3. You want to predict house prices based on square footage. Which approach fits best?}
\begin{enumerate}[label=\alph*)]
    \item K-Means Clustering
    \item Linear Regression
    \item Decision Tree Classification
    \item Image recognition with CNN
\end{enumerate}

\question{Q4. On a scatter plot, what does a good linear regression fit look like?}
\begin{enumerate}[label=\alph*)]
    \item All points stacked vertically
    \item Points scattered randomly with no pattern
    \item Points roughly following a straight line with the regression line passing through them
    \item Points forming a perfect circle
\end{enumerate}

\question{Q5. You have a dataset with 3 features: temperature, humidity, and ``temperature $\times$ 2''. You apply multiple linear regression. What problem will you encounter?}
\begin{enumerate}[label=\alph*)]
    \item The model will train faster
    \item Perfect multicollinearity: the design matrix is singular and the closed-form solution fails
    \item The model will overfit
    \item The MSE will be negative
\end{enumerate}

\question{Q6. Linear Regression assumes a linear relationship between inputs and outputs. Which of these relationships can it still model correctly?}
\begin{enumerate}[label=\alph*)]
    \item $y = \sin(x)$
    \item $y = 3x_1^2 + 2x_2$ (if $x_1^2$ is given as a feature)
    \item $y = e^x$
    \item $y = x_1 \times x_2$ (interaction without feature engineering)
\end{enumerate}

%% ============================================================
\section{The Mathematics (7 Questions)}
%% ============================================================

\question{Q7. In the equation $y = mx + b$, what does ``m'' represent?}
\begin{enumerate}[label=\alph*)]
    \item The y-intercept
    \item The mean of the data
    \item The slope of the line
    \item The number of data points
\end{enumerate}

\question{Q8. If $y = 3x + 2$, what is the predicted value when $x = 4$?}
\begin{enumerate}[label=\alph*)]
    \item 10
    \item 12
    \item 14
    \item 16
\end{enumerate}

\question{Q9. If the slope of a regression line is negative, what does that tell us?}
\begin{enumerate}[label=\alph*)]
    \item There is no relationship between x and y
    \item As x increases, y also increases
    \item As x increases, y decreases
    \item The model has failed
\end{enumerate}

\question{Q10. If $y = -2x + 10$, what is y when $x = 3$?}
\begin{enumerate}[label=\alph*)]
    \item 4
    \item 16
    \item $-4$
    \item 6
\end{enumerate}

\question{Q11. What does the y-intercept (b) represent on a graph?}
\begin{enumerate}[label=\alph*)]
    \item The steepness of the line
    \item Where the line crosses the x-axis
    \item Where the line crosses the y-axis (when $x = 0$)
    \item The average of all y values
\end{enumerate}

\question{Q12. The normal equation for linear regression is $\mathbf{w} = (X^TX)^{-1}X^Ty$. If $X$ is a $100 \times 3$ matrix, what is the shape of $\mathbf{w}$?}
\begin{enumerate}[label=\alph*)]
    \item $100 \times 1$
    \item $3 \times 1$
    \item $3 \times 100$
    \item $1 \times 3$
\end{enumerate}

\question{Q13. Why does the closed-form solution (normal equation) become impractical for very large datasets?}
\begin{enumerate}[label=\alph*)]
    \item It requires more training epochs
    \item Inverting $X^TX$ has $O(n^3)$ complexity in the number of features, which is too slow for high-dimensional data
    \item It always underfits
    \item It cannot handle negative values
\end{enumerate}

%% ============================================================
\section{Cost Function (8 Questions)}
%% ============================================================

\question{Q14. Why do we SQUARE the errors in MSE instead of just summing them?}
\begin{enumerate}[label=\alph*)]
    \item To make computation faster
    \item Because squared numbers are always positive
    \item So positive and negative errors do not cancel out, and large errors are penalized more
    \item It is just a convention with no mathematical reason
\end{enumerate}

\question{Q15. If you drag the regression line far away from the data points, what happens to the MSE?}
\begin{enumerate}[label=\alph*)]
    \item It stays the same
    \item It decreases
    \item It increases dramatically
    \item It becomes negative
\end{enumerate}

\question{Q16. An MSE of 0 means:}
\begin{enumerate}[label=\alph*)]
    \item The model is terrible
    \item The model perfectly predicts every data point
    \item The model has not been trained yet
    \item There is a bug in the code
\end{enumerate}

\question{Q17. Two models have MSE of 25 and 100 respectively. Which is better and by how much?}
\begin{enumerate}[label=\alph*)]
    \item The second model is better because 100 > 25
    \item The first model is 4x better: lower MSE means less error
    \item They are equally good
    \item Cannot be determined from MSE alone
\end{enumerate}

\question{Q18. What is the ``residual'' in regression?}
\begin{enumerate}[label=\alph*)]
    \item The slope of the line
    \item The difference between the actual value and the predicted value
    \item The learning rate
    \item The total number of data points
\end{enumerate}

\question{Q19. MSE uses squared errors. MAE uses absolute errors. When would you prefer MAE over MSE?}
\begin{enumerate}[label=\alph*)]
    \item When you want to penalize large errors more
    \item When your data has many outliers and you want the model to be robust to them
    \item When you need a differentiable loss function
    \item Always: MAE is strictly better
\end{enumerate}

\question{Q20. Given points $(1,2)$, $(2,4)$, $(3,5)$ and the line $y = 1.5x + 0.5$, what is the MSE?}
\begin{enumerate}[label=\alph*)]
    \item 0
    \item 0.0833
    \item 0.25
    \item 0.5
\end{enumerate}

\question{Q21. You compute the gradient of MSE and find $\frac{\partial L}{\partial w} = 0$ and $\frac{\partial L}{\partial b} = 0$. What does this tell you?}
\begin{enumerate}[label=\alph*)]
    \item The model has not started training yet
    \item The current $w$ and $b$ are at a critical point (minimum) of the loss function
    \item The learning rate is zero
    \item The data has no pattern
\end{enumerate}

%% ============================================================
\section{Gradient Descent (8 Questions)}
%% ============================================================

\question{Q22. What happens if the learning rate is TOO HIGH?}
\begin{enumerate}[label=\alph*)]
    \item The model converges faster
    \item The model overshoots the minimum and may never converge
    \item The model underfits
    \item Nothing changes
\end{enumerate}

\question{Q23. You are training a model and the loss keeps bouncing up and down without decreasing. What is the most likely cause?}
\begin{enumerate}[label=\alph*)]
    \item The dataset is too small
    \item The learning rate is too high
    \item The model needs more features
    \item The bias term is wrong
\end{enumerate}

\question{Q24. In Gradient Descent, the gradient tells us:}
\begin{enumerate}[label=\alph*)]
    \item The exact location of the minimum
    \item The direction of steepest INCREASE of the loss
    \item The learning rate to use
    \item How many epochs are needed
\end{enumerate}

\question{Q25. Your model is training very slowly, taking thousands of epochs. The loss decreases but extremely gradually. What should you do?}
\begin{enumerate}[label=\alph*)]
    \item Decrease the learning rate
    \item Increase the learning rate
    \item Add more data points
    \item Remove the bias term
\end{enumerate}

\question{Q26. On the cost function surface, what does the lowest point represent?}
\begin{enumerate}[label=\alph*)]
    \item The worst possible model
    \item The starting point of training
    \item The optimal weights that minimize error
    \item The learning rate value
\end{enumerate}

\question{Q27. What is the partial derivative of MSE $= \frac{1}{n}\sum(wx_i + b - y_i)^2$ with respect to $w$?}
\begin{enumerate}[label=\alph*)]
    \item $\frac{2}{n}\sum(wx_i + b - y_i)$
    \item $\frac{2}{n}\sum(wx_i + b - y_i) \cdot x_i$
    \item $\frac{1}{n}\sum x_i$
    \item $\frac{2}{n}\sum y_i$
\end{enumerate}

\question{Q28. You plot the loss curve and notice it decreases, then increases, then decreases again in a zigzag pattern. What is the best fix?}
\begin{enumerate}[label=\alph*)]
    \item Use a larger learning rate
    \item Use learning rate decay: start high and gradually reduce it
    \item Remove all features
    \item Use MSE instead of MAE
\end{enumerate}

\question{Q29. For a single data point $(x=2, y=6)$ with current $w=1$ and $b=0$, what is $\partial\text{MSE}/\partial w$?}
\begin{enumerate}[label=\alph*)]
    \item $-16$
    \item $-8$
    \item $8$
    \item $16$
\end{enumerate}

%% ============================================================
\section{Training Process (6 Questions)}
%% ============================================================

\question{Q30. What is one ``epoch'' in the training process?}
\begin{enumerate}[label=\alph*)]
    \item One update to the weights
    \item One complete pass through the entire dataset
    \item One prediction
    \item One gradient calculation
\end{enumerate}

\question{Q31. After 10 epochs the loss is 500. After 100 epochs it is 50. After 200 epochs it is 49.8. What should you do?}
\begin{enumerate}[label=\alph*)]
    \item Keep training for 1000 more epochs
    \item Stop training: the model has converged
    \item Increase the learning rate dramatically
    \item Start over with random weights
\end{enumerate}

\question{Q32. What happens to the weights during each training step?}
\begin{enumerate}[label=\alph*)]
    \item They are randomly reassigned
    \item They are updated in the direction that reduces the loss
    \item They stay the same until the final epoch
    \item They are always increased
\end{enumerate}

\question{Q33. If weight $= 2.0$, gradient $= 0.5$, learning rate $= 0.1$, what is the new weight?}
\begin{enumerate}[label=\alph*)]
    \item 2.05
    \item 1.95
    \item 2.5
    \item 1.5
\end{enumerate}

\question{Q34. What is the difference between batch, mini-batch, and stochastic gradient descent?}
\begin{enumerate}[label=\alph*)]
    \item They use different loss functions
    \item Batch uses all data per update, mini-batch uses a subset, SGD uses a single sample
    \item SGD is always faster and more accurate
    \item There is no practical difference
\end{enumerate}

\question{Q35. A model trained for 50 epochs has training MSE=10 and test MSE=100. What is happening?}
\begin{enumerate}[label=\alph*)]
    \item The model is underfitting
    \item The model is overfitting: it memorized training data but does not generalize
    \item The model is perfectly trained
    \item The test set is too large
\end{enumerate}

%% ============================================================
\section{Failure Cases (7 Questions)}
%% ============================================================

\question{Q36. A dataset contains a clear quadratic (curved) relationship. You apply Linear Regression. What happens?}
\begin{enumerate}[label=\alph*)]
    \item It fits perfectly
    \item It underfits: the straight line cannot capture the curve
    \item It overfits
    \item The MSE becomes zero
\end{enumerate}

\question{Q37. A dataset has one extreme outlier. What effect does it have on the regression line?}
\begin{enumerate}[label=\alph*)]
    \item No effect at all
    \item It pulls the line toward the outlier, distorting predictions for other points
    \item It improves the model accuracy
    \item It only affects the intercept
\end{enumerate}

\question{Q38. Adding heavy random noise to a dataset will likely cause:}
\begin{enumerate}[label=\alph*)]
    \item The MSE to decrease
    \item Perfect predictions
    \item Higher MSE and less reliable predictions
    \item The model to overfit
\end{enumerate}

\question{Q39. What is ``multicollinearity'' and why is it a problem?}
\begin{enumerate}[label=\alph*)]
    \item When features are unrelated: makes the model too simple
    \item When two or more features are highly correlated: makes weight estimates unstable
    \item When there is only one feature: limits accuracy
    \item When the target variable has multiple values
\end{enumerate}

\question{Q40. If your data forms a U-shape (parabola), what will the regression line look like?}
\begin{enumerate}[label=\alph*)]
    \item A perfect curve matching the data
    \item A horizontal line cutting through the middle, missing both ends
    \item A vertical line
    \item Two separate lines
\end{enumerate}

\question{Q41. Your linear regression has $R^2 = 0.15$. What does this mean?}
\begin{enumerate}[label=\alph*)]
    \item The model explains 85\% of variance: it is great
    \item The model explains only 15\% of variance: the relationship is likely nonlinear or key features are missing
    \item $R^2$ of 0.15 is always acceptable
    \item The model has converged optimally
\end{enumerate}

\question{Q42. What is heteroscedasticity?}
\begin{enumerate}[label=\alph*)]
    \item When the mean of errors is non-zero
    \item When the variance of residuals changes across different values of x, violating the constant-variance assumption
    \item When errors are correlated with each other
    \item When the data has too few points
\end{enumerate}

%% ============================================================
\section{Comparison (7 Questions)}
%% ============================================================

\question{Q43. What is the key advantage of Linear Regression over Neural Networks?}
\begin{enumerate}[label=\alph*)]
    \item It is more accurate
    \item It can handle nonlinear data
    \item It is simpler, faster, and more interpretable
    \item It requires more data
\end{enumerate}

\question{Q44. Which method should you try FIRST for a dataset with a clear linear trend?}
\begin{enumerate}[label=\alph*)]
    \item Deep Neural Network
    \item Random Forest
    \item Linear Regression
    \item Convolutional Neural Network
\end{enumerate}

\question{Q45. You have 50 data points with a clear linear trend. A colleague suggests a neural network with 10 layers. What do you recommend?}
\begin{enumerate}[label=\alph*)]
    \item Agree: deeper is always better
    \item Use Linear Regression: it is sufficient and will not overfit
    \item Use at least 20 layers
    \item Do not use any model
\end{enumerate}

\question{Q46. What does ``interpretability'' mean in the context of models?}
\begin{enumerate}[label=\alph*)]
    \item How fast the model trains
    \item How easily humans can understand WHY the model makes its predictions
    \item How many parameters the model has
    \item The programming language used
\end{enumerate}

\question{Q47. Your image recognition task needs to classify cats vs dogs. Should you use Linear Regression?}
\begin{enumerate}[label=\alph*)]
    \item Yes: it works for everything
    \item No: this is a classification task with complex, nonlinear patterns requiring a different approach
    \item Yes, but only with a high learning rate
    \item Only if you have less than 100 images
\end{enumerate}

\question{Q48. What is regularization (L1/L2)?}
\begin{enumerate}[label=\alph*)]
    \item It increases the learning rate for faster training
    \item It adds a penalty term to the loss to prevent weights from becoming too large, reducing overfitting
    \item It removes outliers from the dataset
    \item It converts regression into classification
\end{enumerate}

\question{Q49. You apply Ridge Regression (L2) with $\lambda=1000$. What happens to the weights?}
\begin{enumerate}[label=\alph*)]
    \item They become very large
    \item They are pushed very close to zero, severely underfitting the data
    \item They are unaffected
    \item They become negative
\end{enumerate}

%% ============================================================
\section{Multivariate Regression (13 Questions)}
%% ============================================================

\question{Q50. What does multivariable linear regression extend compared to simple linear regression?}
\begin{enumerate}[label=\alph*)]
    \item The number of output variables
    \item The number of input features
    \item The number of training epochs
    \item The number of loss functions
\end{enumerate}

\question{Q51. In the equation $\hat{y} = w_1x_1 + w_2x_2 + w_3x_3 + b$, how many features does the model use?}
\begin{enumerate}[label=\alph*)]
    \item 1
    \item 2
    \item 3
    \item 4
\end{enumerate}

\question{Q52. You are predicting house prices using size, bedrooms, and age. Which is the correct model form?}
\begin{enumerate}[label=\alph*)]
    \item $\hat{y} = w \cdot \text{size} + b$
    \item $\hat{y} = w_1 \cdot \text{size} + w_2 \cdot \text{bedrooms} + w_3 \cdot \text{age} + b$
    \item $\hat{y} = w \cdot (\text{size} + \text{bedrooms} + \text{age}) + b$
    \item $\hat{y} = w_1 \cdot \text{size} \times w_2 \cdot \text{bedrooms} \times w_3 \cdot \text{age} + b$
\end{enumerate}

\question{Q53. In the matrix formulation $\hat{Y} = XW + b$, what are the dimensions of $X$ if there are 100 samples and 5 features?}
\begin{enumerate}[label=\alph*)]
    \item $5 \times 100$
    \item $100 \times 5$
    \item $100 \times 100$
    \item $5 \times 5$
\end{enumerate}

\question{Q54. What is the Normal Equation for finding optimal weights?}
\begin{enumerate}[label=\alph*)]
    \item $W^* = X^Ty$
    \item $W^* = (X^TX)^{-1}X^Ty$
    \item $W^* = X^{-1}y$
    \item $W^* = (XX^T)^{-1}Xy$
\end{enumerate}

\question{Q55. Feature ``size'' ranges from 50--500 while ``bedrooms'' ranges from 1--6. What problem does this cause?}
\begin{enumerate}[label=\alph*)]
    \item The model cannot converge
    \item Gradient descent zig-zags and converges slowly
    \item The bias becomes negative
    \item The model always overfits
\end{enumerate}

\question{Q56. What does Min-Max Normalization scale features to?}
\begin{enumerate}[label=\alph*)]
    \item $[-1, 1]$
    \item $[0, 1]$
    \item Mean$=0$, Std$=1$
    \item $[0, 100]$
\end{enumerate}

\question{Q57. What does Z-score standardization produce?}
\begin{enumerate}[label=\alph*)]
    \item Values between 0 and 1
    \item Values between $-1$ and 1
    \item Mean $= 0$ and standard deviation $= 1$
    \item All positive values
\end{enumerate}

\question{Q58. What is feature engineering?}
\begin{enumerate}[label=\alph*)]
    \item Removing all features from the dataset
    \item Creating new features from existing ones to improve predictions
    \item Training a model without any features
    \item Using only one feature at a time
\end{enumerate}

\question{Q59. You have 50,000 features. Should you use the Normal Equation or Gradient Descent?}
\begin{enumerate}[label=\alph*)]
    \item Normal Equation: it is always exact
    \item Gradient Descent: Normal Equation is $O(n^3)$ and too slow for many features
    \item Either works equally well
    \item Neither: use a neural network instead
\end{enumerate}

\question{Q60. What is multicollinearity and why is it problematic?}
\begin{enumerate}[label=\alph*)]
    \item When features are uncorrelated with the target
    \item When features are highly correlated with each other, causing unstable weights
    \item When the model has too many layers
    \item When the learning rate is too high
\end{enumerate}

\question{Q61. Given $\hat{y} = 3x_1 + 2x_2 - x_3 + 5$, what is the prediction for $x_1=2, x_2=4, x_3=1$?}
\begin{enumerate}[label=\alph*)]
    \item 16
    \item 18
    \item 14
    \item 20
\end{enumerate}

\question{Q62. Your features have very different scales: $x_1 \in [0,1]$ and $x_2 \in [0, 1000000]$. You run gradient descent without normalization. What happens?}
\begin{enumerate}[label=\alph*)]
    \item Nothing: gradient descent handles all scales equally
    \item The loss surface becomes elongated, causing slow zigzag convergence
    \item The model only learns $x_2$ and ignores $x_1$
    \item The MSE becomes negative
\end{enumerate}

%% ============================================================
\section{Data Intuition (5 Questions)}
%% ============================================================

\question{Q63. What does a Pearson correlation coefficient of $r = 0$ mean?}
\begin{enumerate}[label=\alph*)]
    \item Perfect positive relationship
    \item No linear relationship
    \item Perfect negative relationship
    \item The data has errors
\end{enumerate}

\question{Q64. If data points form a tight cluster along a downward-sloping line, what is the likely correlation?}
\begin{enumerate}[label=\alph*)]
    \item $r \approx +0.9$
    \item $r \approx -0.9$
    \item $r \approx 0$
    \item $r \approx +0.5$
\end{enumerate}

\question{Q65. Ice cream sales and drowning deaths have $r \approx 0.85$. Does this mean ice cream causes drowning?}
\begin{enumerate}[label=\alph*)]
    \item Yes, the high correlation proves causation
    \item No: both are caused by a confounding variable (hot weather)
    \item Only if $r > 0.95$
    \item This correlation is impossible
\end{enumerate}

\question{Q66. Which scenario would produce a Pearson $r \approx 0$ despite a clear pattern in the data?}
\begin{enumerate}[label=\alph*)]
    \item A perfectly linear positive trend
    \item A U-shaped (quadratic) relationship
    \item A moderate negative trend
    \item A strong positive trend with noise
\end{enumerate}

\question{Q67. If $r = -1.0$, what does the scatter plot look like?}
\begin{enumerate}[label=\alph*)]
    \item Random scatter
    \item All points exactly on an upward line
    \item All points exactly on a downward line
    \item A circular pattern
\end{enumerate}

%% ============================================================
\section{Bias-Variance Tradeoff (5 Questions)}
%% ============================================================

\question{Q68. What is ``underfitting'' in machine learning?}
\begin{enumerate}[label=\alph*)]
    \item The model memorizes training data
    \item The model is too simple to capture the underlying pattern
    \item The model works perfectly on test data
    \item The model has too many parameters
\end{enumerate}

\question{Q69. What happens when a model overfits?}
\begin{enumerate}[label=\alph*)]
    \item It performs well on both train and test data
    \item It has high training error
    \item It memorizes training data but fails on new data
    \item It is too simple
\end{enumerate}

\question{Q70. Your model has train MSE $= 0.5$ and test MSE $= 25$. What is the diagnosis?}
\begin{enumerate}[label=\alph*)]
    \item Underfitting
    \item Good fit
    \item Overfitting
    \item Insufficient data
\end{enumerate}

\question{Q71. In the bias-variance tradeoff, increasing model complexity typically:}
\begin{enumerate}[label=\alph*)]
    \item Increases bias and decreases variance
    \item Decreases bias and increases variance
    \item Decreases both bias and variance
    \item Increases both bias and variance
\end{enumerate}

\question{Q72. A degree-15 polynomial is fit to 20 data points. What will likely happen?}
\begin{enumerate}[label=\alph*)]
    \item Perfect generalization
    \item High bias, low variance
    \item Low bias, high variance (overfitting)
    \item The model cannot be trained
\end{enumerate}

%% ============================================================
\section{Evaluation Metrics (5 Questions)}
%% ============================================================

\question{Q73. What does $R^2 = 1.0$ mean?}
\begin{enumerate}[label=\alph*)]
    \item The model has 100\% error
    \item The model explains all variance in the target
    \item The model is overfitting
    \item The features are uncorrelated
\end{enumerate}

\question{Q74. Which metric is MOST sensitive to outliers?}
\begin{enumerate}[label=\alph*)]
    \item MAE (Mean Absolute Error)
    \item MSE (Mean Squared Error)
    \item Median Error
    \item $R^2$
\end{enumerate}

\question{Q75. Predictions: $[10, 12, 15]$. Actual: $[11, 10, 14]$. What is the MAE?}
\begin{enumerate}[label=\alph*)]
    \item 1.0
    \item 1.33
    \item 2.0
    \item 1.67
\end{enumerate}

\question{Q76. Your model has RMSE $= 5.0$ predicting house prices in thousands. What does this mean practically?}
\begin{enumerate}[label=\alph*)]
    \item Predictions are off by \$5 on average
    \item Predictions are off by about \$5,000 on average
    \item The model explains 5\% of variance
    \item The model has 5 features
\end{enumerate}

\question{Q77. Can $R^2$ be negative? If so, when?}
\begin{enumerate}[label=\alph*)]
    \item No, $R^2$ is always between 0 and 1
    \item Yes: when the model is worse than predicting the mean
    \item Yes: only for nonlinear models
    \item Yes: when there are too many features
\end{enumerate}

%% ============================================================
\section{Linear Regression to Neural Networks (5 Questions)}
%% ============================================================

\question{Q78. A single neuron with no activation function is equivalent to:}
\begin{enumerate}[label=\alph*)]
    \item A decision tree
    \item Linear regression
    \item A random forest
    \item Logistic regression
\end{enumerate}

\question{Q79. Why do neural networks need activation functions?}
\begin{enumerate}[label=\alph*)]
    \item To make training faster
    \item To introduce non-linearity so the network can model complex patterns
    \item To reduce the number of parameters
    \item To prevent overfitting
\end{enumerate}

\question{Q80. What does the ReLU activation function do?}
\begin{enumerate}[label=\alph*)]
    \item Outputs the input if positive, 0 otherwise
    \item Squishes values between 0 and 1
    \item Squares the input
    \item Normalizes the output
\end{enumerate}

\question{Q81. If you stack 3 linear layers without activations ($y = W_3(W_2(W_1x + b_1) + b_2) + b_3$), what do you get?}
\begin{enumerate}[label=\alph*)]
    \item A deep neural network
    \item A single linear transformation (still linear regression)
    \item A cubic function
    \item A model that cannot be trained
\end{enumerate}

\question{Q82. A neuron has weights $[2, -1]$, bias $= 3$, and ReLU activation. Input $= [1, 5]$. What is the output?}
\begin{enumerate}[label=\alph*)]
    \item 0
    \item 6
    \item $-2$
    \item 0 (because ReLU clips negative values)
\end{enumerate}

%% ============================================================
%%                      ANSWER KEY
%% ============================================================
\newpage
\section{Answer Key}

All correct answers with brief explanations.

\vspace{0.5em}

\renewcommand{\arraystretch}{1.25}

\begin{longtable}{|c|c|p{9.5cm}|}
\hline
\textbf{Q\#} & \textbf{Ans} & \textbf{Explanation} \\
\hline
\endfirsthead
\hline
\textbf{Q\#} & \textbf{Ans} & \textbf{Explanation} \\
\hline
\endhead
\hline
\endfoot
1 & b & Linear Regression finds the best straight line ($y = wx + b$) that minimizes the distance between predictions and actual values. \\
\hline
2 & b & It predicts continuous numeric values (e.g., price, temperature) rather than discrete categories. \\
\hline
3 & b & Predicting a continuous value (price) from a numeric feature (square footage) is a classic Linear Regression problem. \\
\hline
4 & c & A good fit shows points roughly following a linear trend, with the regression line passing through the center. \\
\hline
5 & b & The third feature is an exact linear combination of the first, making $(X^TX)$ singular. \\
\hline
6 & b & Linear regression is linear in its parameters, not its features. If $x_1^2$ is provided as a feature, the model fits it perfectly. \\
\hline
7 & c & The slope (m) determines how steep the line is. \\
\hline
8 & c & $y = 3(4) + 2 = 14$ \\
\hline
9 & c & A negative slope means as x increases, y decreases. \\
\hline
10 & a & $y = -2(3) + 10 = 4$ \\
\hline
11 & c & The y-intercept is the value of y when $x = 0$. \\
\hline
12 & b & $X^TX$ is $3\times3$, so $\mathbf{w} = (3\times3)(3\times1) = 3\times1$. \\
\hline
13 & b & Inverting $X^TX$ has $O(p^3)$ complexity where p is the number of features. \\
\hline
14 & c & Squaring prevents cancellation of positive/negative errors and penalizes larger errors more. \\
\hline
15 & c & Moving the line away increases squared errors dramatically. \\
\hline
16 & b & MSE $= 0$ means every prediction exactly matches the actual value. \\
\hline
17 & b & Lower MSE means less error. The first model has 4x less error. \\
\hline
18 & b & Residual $=$ actual $-$ predicted. \\
\hline
19 & b & MAE is more robust to outliers since it does not square errors. \\
\hline
20 & b & Predictions: 2, 3.5, 5. Errors$^2$: 0, 0.25, 0. MSE $= 0.25/3 = 0.0833$. \\
\hline
21 & b & Zero gradients mean you are at the global minimum (MSE is convex for linear regression). \\
\hline
22 & b & Too-high learning rate causes overshooting and divergence. \\
\hline
23 & b & Oscillating loss typically indicates a learning rate that is too large. \\
\hline
24 & b & The gradient points in the direction of steepest increase; we move opposite to decrease loss. \\
\hline
25 & b & Very slow convergence means the learning rate is too small. \\
\hline
26 & c & The lowest point is the optimal weights that minimize error. \\
\hline
27 & b & By chain rule: $\frac{2}{n}\sum(wx_i + b - y_i) \cdot x_i$. \\
\hline
28 & b & Learning rate decay starts high for fast progress, then reduces for stable convergence. \\
\hline
29 & a & $\hat{y}=2$, error $= -4$. $\partial\text{MSE}/\partial w = 2(-4)(2) = -16$. \\
\hline
30 & b & An epoch is one complete pass through all training data. \\
\hline
31 & b & Loss barely changed ($50 \to 49.8$): the model has converged. \\
\hline
32 & b & Weights are updated in the direction that reduces the loss. \\
\hline
33 & b & $2.0 - 0.1 \times 0.5 = 1.95$ \\
\hline
34 & b & Batch uses all data, mini-batch a subset, SGD one sample per update. \\
\hline
35 & b & Large gap between train and test error is overfitting. \\
\hline
36 & b & A straight line cannot capture a curve: underfitting. \\
\hline
37 & b & MSE penalizes large errors, pulling the line toward the outlier. \\
\hline
38 & c & Noise obscures the true pattern, raising MSE. \\
\hline
39 & b & Highly correlated features make weight estimates unstable. \\
\hline
40 & b & A straight line cuts through the middle of a U-shape, missing both ends. \\
\hline
41 & b & Only 15\% of variance is explained; likely nonlinear or missing features. \\
\hline
42 & b & Error variance changes with x, violating the constant-variance assumption. \\
\hline
43 & c & Linear Regression is simpler, faster, and more interpretable. \\
\hline
44 & c & Always start simple with Linear Regression for linear data. \\
\hline
45 & b & For 50 points with a linear trend, LR is sufficient; a deep network would overfit. \\
\hline
46 & b & How easily humans can understand why the model makes its predictions. \\
\hline
47 & b & Image classification needs a CNN, not regression. \\
\hline
48 & b & Regularization penalizes large weights to reduce overfitting. \\
\hline
49 & b & Very high $\lambda$ pushes all weights toward zero, causing underfitting. \\
\hline
50 & b & Extends from one input feature to multiple input features. \\
\hline
51 & c & Three features: $x_1, x_2, x_3$. b is the bias. \\
\hline
52 & b & Each feature gets its own independent weight. \\
\hline
53 & b & $X$ is (samples $\times$ features) $= 100 \times 5$. \\
\hline
54 & b & $W^* = (X^TX)^{-1}X^Ty$ \\
\hline
55 & b & Different scales cause elongated contours and slow zig-zag convergence. \\
\hline
56 & b & Maps values to $[0, 1]$ using $(x - \min)/(\max - \min)$. \\
\hline
57 & c & Centers at mean $= 0$ with standard deviation $= 1$. \\
\hline
58 & b & Creating new features from existing ones (e.g., $x^2$, price/size). \\
\hline
59 & b & Normal Equation is $O(n^3)$; gradient descent scales better. \\
\hline
60 & b & Correlated features cause unstable weight estimates. \\
\hline
61 & b & $3(2) + 2(4) - 1(1) + 5 = 18$ \\
\hline
62 & b & Unscaled features create an elongated loss surface. \\
\hline
63 & b & No linear relationship between the variables. \\
\hline
64 & b & Tight downward cluster $= $ strong negative correlation. \\
\hline
65 & b & Correlation $\neq$ causation; hot weather is the confounding variable. \\
\hline
66 & b & U-shaped patterns cancel out, giving $r \approx 0$ despite a clear pattern. \\
\hline
67 & c & $r = -1.0$: all points on a downward line. \\
\hline
68 & b & The model is too simple to capture the pattern. \\
\hline
69 & c & Memorizes training data but fails on new data. \\
\hline
70 & c & Huge train/test gap $=$ overfitting. \\
\hline
71 & b & More complexity $=$ less bias, more variance. \\
\hline
72 & c & 16 parameters for 20 points $=$ extreme overfitting. \\
\hline
73 & b & The model explains all variance in the target. \\
\hline
74 & b & MSE squares errors, so outliers dominate. \\
\hline
75 & b & $(1+2+1)/3 \approx 1.33$ \\
\hline
76 & b & RMSE is in target units: \$5,000 off on average. \\
\hline
77 & b & $R^2 < 0$ when the model is worse than predicting the mean. \\
\hline
78 & b & A neuron with no activation $= $ linear regression. \\
\hline
79 & b & Activations introduce non-linearity for complex patterns. \\
\hline
80 & a & ReLU$(x) = \max(0, x)$. \\
\hline
81 & b & Stacked linear layers $=$ one linear transformation. \\
\hline
82 & a & $z = 2(1) + (-1)(5) + 3 = 0$. ReLU$(0) = 0$. \\
\hline
\end{longtable}

\end{document}
